\section{Safety: what stops it before it hurts someone?}
\label{sec:safety}

The clinician's duty of non-maleficence makes safety the second question. It is not
enough that a system usually behaves; the clinician needs to know what concretely
stops an irreversible action before it can reach the patient. The fear behind this
question is not abstract. The landmark report \textit{To Err Is Human} estimated
that tens of thousands of Americans die each year from preventable medical error
\cite{kohn2000toerr}, and a later analysis ranked medical error among the leading
causes of death in the United States \cite{makary2016medical}. A safety case that
deserves a clinician's trust must therefore name the specific thing that prevents
the worst outcome, not merely promise that it is unlikely.

In this framework that specific thing is the ten-gate VVUQ examination. Every
candidate action is generated under verification before generation, in which
Claude Code proposes and Codex reviews independently, and then must pass ten gates
before it may proceed. The examination resolves to ACCEPT (act under audit), ESCALATE (hand the
decision to a qualified human), or BLOCK (stop the action before it reaches the
patient). The gates are not advisory. They are encoded in a standard operating
procedure, the written, version-controlled document that fixes who does what, in
what order, and under which thresholds, so that the same action is checked the same
way every time and no step depends on an individual's memory at the moment of
decision.

The gates fall into three tiers, and \autoref{tab:safety} sets out what each tier
checks and what happens when a check fails. The first four gates test numeric
thresholds (accuracy, calibration, uncertainty, and tolerance bounds), so that a
quantitatively deficient action never advances. Gates five through nine bind the
action to external standards and protocol constraints, so that compliance is
verified rather than assumed. Gate ten is the catastrophe predicate: if the
contemplated action matches a predicate for irreversible harm, the system issues a
BLOCK before the patient, regardless of how confident the earlier gates were. This
is the hard stop the question demands, and Figure 4 shows it as the final fork
between contain and BLOCK.

Safety also depends on what happens after a stop. A BLOCK or an ESCALATE is only
useful if a competent human is actually there to receive it, at any hour. The
author's adverse-event lineage describes a threefold humanoid response team
configured for 24/7 coverage, so that an escalated case meets a qualified responder
rather than a queue \cite{Kawchak2026ThreefoldHumanoid247}. The illustrative
H. R. 9510 would make this coverage and the gate suite part of the binding trial
record rather than a vendor's promise \cite{Kawchak2026HR9510Bill}, so that the
existence of the hard stop, and of the team behind it, can be confirmed by an
outside party.

\begin{cfwfig}
\node[cfone] (a) at (0,0) {Candidate action}; \node[cfdec] (b) at (3.6,0) {Gate 10};
\node[cfhope] (c) at (7.0,0.8) {Contain}; \node[cfrisk] (d) at (7.0,-0.8) {BLOCK};
\draw[cfedge] (a)--(b); \draw[cfedge] (b)--(c); \draw[cfedge] (b)--(d);
\end{cfwfig}
\vspace{-0.2cm}
\cfwcaption{Figure 4. The catastrophe-predicate gate (Mermaid 04 and 07 reproduced
as colored TikZ): a candidate action that matches the irreversible-harm predicate
is contained or issued a BLOCK before it can reach the patient.}

\needspace{9\baselineskip}\begin{table}[H]
\footnotesize
\begin{tabularx}{\textwidth}{@{}L{3cm} Y L{3.4cm}@{}}
\toprule
\textbf{Gate} & \textbf{What it checks} & \textbf{If it fails}\\
\midrule
Gates 1 to 4 (numeric thresholds) & Accuracy, calibration, uncertainty, and tolerance bounds against pre-set numeric limits & The action is held and ESCALATED to a qualified human for review\\
Gates 5 to 9 (standards binding) & Conformance to external standards, protocol constraints, and the signed standard operating procedure & The deviation is logged and the action is blocked pending correction\\
Gate 10 (catastrophe predicate) & Whether the action matches a predicate for irreversible or catastrophic harm & The system issues a BLOCK before the patient, overriding earlier passes\\
\bottomrule
\end{tabularx}
\caption{The ten gates in three tiers: what each checks and the consequence of a
failure, ending in the catastrophe-predicate BLOCK.}
\label{tab:safety}
\end{table}

A framework cannot promise that nothing will go wrong. It can promise a hard stop
before the worst outcomes, and a record of every action that does proceed. The
gate, the standard operating procedure behind it, and the 24/7 team that answers an
escalation are the concrete answer to the question every clinician carries to the
bedside: would I let this system near my own family member, knowing exactly what
stops it before it hurts someone?
